\clearpage
\SourceLeaf{101}%

\[
\boldsymbol{P}(0)=-1,\qquad \boldsymbol{P}\varphi=\boldsymbol{P}(-\varphi),\qquad
\boldsymbol{P}\psi=0,\qquad \boldsymbol{P}(2\varphi)=2,\qquad \boldsymbol{P}(3\psi)=\infty,
\]
\[
\boldsymbol{P}(\varphi+3\psi)=\frac{2-\boldsymbol{P}(\varphi)}{1+\boldsymbol{P}(\varphi)},
\qquad \boldsymbol{P}(1\tfrac12\psi)=\sqrt{3}-1.
\]

\noindent Wenn
\[
\boldsymbol{P}\varphi=s,\qquad \boldsymbol{P}(\varphi+2\psi)=c,
\]
\noindent so ist
\[
s\,s\,c\,c=4s+4c,\qquad c=\frac{2(1+\sqrt{1+s^3})}{s\,s},
\]
\[
\boldsymbol{P}(3\psi+\varphi+\varphi')=
\frac{2+s s'(s+s')-\tfrac12(c s s-2)(c' s' s'-2)}{(s-s')^2},
\]
\[
\boldsymbol{P}(3\psi+\varphi-\varphi')=
\frac{2+s s'(s+s')+\tfrac12(c s s-2)(c' s' s'-2)}{(s-s')^2}.
\]



\Needspace{10\baselineskip}
Die vorigen Sätze gelten für die angezeigte Bedeutung unserer Begriffe. Allein lassen wir
$\Pi z$ den Werth des Integrals $\int \frac{dx}{\BareRoot{(1+x^3)}}$ von $x=-1$ bis $x=z-1$
bedeuten (welches sinnlich gemacht wird, wenn man die Abscissenlinie um 1 herunterrückt),
so wird:

\[
\boldsymbol{P}(0)=0,\quad \boldsymbol{P}\psi=1,\quad \boldsymbol{P}(2\psi)=3,\quad
\boldsymbol{P}(3\psi)=\infty,\quad \boldsymbol{P}(3\psi-\varphi)=\frac{3}{\boldsymbol{P}(\psi)},
\]
\[
\boldsymbol{P}(2\varphi)=
\frac{12\boldsymbol{P}\varphi\left(\boldsymbol{P}(\varphi)^2-3\boldsymbol{P}\varphi+3\right)}
{\left(\boldsymbol{P}(\varphi)^2-3\right)^2},
\]
\[
\boldsymbol{P}(\varphi+\varphi')=
\frac{\begin{gathered}3(s+s')s s'-18 s s'+9(s+s')\\{}+6\sqrt{(s^3-3ss+3s)(s'^3-3s's'+3s')}\end{gathered}}{(ss'-3)^2}.
\]



\Needspace{5\baselineskip}
\[
\Pi(1+z)=\psi+z-\frac14\cdot\frac12 z^4+\frac17\cdot\frac{1.3}{2.4}z^7
-\frac{1}{10}\cdot\frac{1.3.5}{2.4.6}z^{10}+\cdots,
\]
\[
\Pi(1+z)=3\psi-2\sqrt{z}\left(\frac1z-\frac17\cdot\frac12\frac{1}{z^4}
+\frac{1}{13}\cdot\frac{1.3}{2.4}\cdot\frac{1}{z^7}-\cdots\right),
\]
\[
\psi=1{,}4021822
\]
\[
\log\psi=0{,}1468045
\]
\[
\log\mathrm{hyp}\,\psi=0{,}3380399.
\]


